% Generated draft: requires mathematical and source review.
\documentclass[11pt]{article}
\usepackage[margin=25mm]{geometry}
\usepackage{fontspec}
\setmainfont{Latin Modern Roman}
\usepackage{amsmath,amssymb,mathtools,unicode-math}
\setmathfont{Latin Modern Math}
\usepackage{enumitem,needspace,xurl,hyperref}
\setlength{\parindent}{0pt}
\setlength{\parskip}{5pt}
\setlength{\emergencystretch}{3em}
\newcommand{\N}{\mathbb{N}}
\newcommand{\Z}{\mathbb{Z}}
\newcommand{\Q}{\mathbb{Q}}
\newcommand{\R}{\mathbb{R}}
\newcommand{\C}{\mathbb{C}}
\newcommand{\F}{\mathbb{F}}
\newcommand{\A}{\mathbb{A}}
\newcommand{\PP}{\mathbb{P}}
\newcommand{\E}{\mathbb{E}}
\newcommand{\eps}{\varepsilon}
\newcommand{\dd}{\,\mathrm{d}}
\begin{document}
\section*{UnsolvedMath ID: 30000944}
\textbf{Source problem code:} catalogue entry
\subsection{Definitions and mathematical statement}
For a complex polynomial \(f\) of degree at least two, define its filled Julia set and Julia set by
\[
 K(f)=\{z\in\C:\sup_{m\geq0}|f^{\circ m}(z)|<\infty\},
 \qquad J(f)=\partial K(f).
\]
These are nonempty compact subsets of \(\C\). Write
\[
 f_j(z)=\sum_{r=0}^d a_{r,j}z^r,\quad a_{d,j}\ne0,\qquad
 a_{r,j}\to a_r,\quad f(z)=\sum_{r=0}^d a_rz^r.
\]
The source fixes the degree \(d\) of \(f_j\), but does not explicitly require \(d\geq2\) or \(a_d\ne0\) for the limit. These conditions are necessary to use the usual nondegenerate compact-set formulation without further conventions.

For nonempty compact \(A,B\subseteq\C\), put
\[
 d^+(A,B)=\sup_{a\in A}\inf_{b\in B}|a-b|,
 \qquad d_H(A,B)=\max\{d^+(A,B),d^+(B,A)\}.
\]
Upper semicontinuity of \(f\mapsto K(f)\) in the nondegenerate degree-\(d\) parameter space means
\[
 d^+(K(f_j),K(f))\longrightarrow0,
\]
while lower semicontinuity of \(f\mapsto J(f)\) means
\[
 d^+(J(f),J(f_j))\longrightarrow0.
\]
Hausdorff continuity requires both directional distances to tend to zero for the relevant set-valued map. The question asks which of these or other compact-set convergence statements hold and under which dynamical hypotheses, rather than assuming unrestricted Hausdorff convergence.

If degree drop is allowed, a different formulation is required. For example, \(f_j(z)=z^2/j\to0\) coefficientwise has \(K(f_j)=\{|z|\leq j\}\) and \(J(f_j)=\{|z|=j\}\), while the bounded-orbit set of the constant limit is all of \(\C\). Thus ordinary nonempty-compact Hausdorff convergence in \(\C\) no longer has the same target. A spherical or local set topology, or exclusion of degree drop, cannot be silently chosen.

The supplied assessment dated 2026-08-21 attributes to Douady upper semicontinuity for filled Julia sets, lower semicontinuity for Julia sets, and continuity under a precise dynamical condition, while citing parabolic implosion as an obstruction to unrestricted continuity. The condition and theorem scope are not reproduced; in particular this summary is not extended to the degree-drop situation.
\subsection{Short English statement}
The question asks how Julia and filled Julia sets vary under coefficient convergence, distinguishing one-sided semicontinuity from Hausdorff continuity. The usual fixed-degree theory requires the limit to retain degree at least two; the source wording leaves that point open, and degree drop changes the compact-set setting.
\subsection{Sources}
\begin{itemize}
\item Local UnsolvedMath registry, permanent ID \(30000944\), exact source code \texttt{catalogue entry}, proposed year 2008; no individual proposer is supplied.
\item \emph{Arbeitsgemeinschaft: Julia Sets of Positive Measure (2009)}, Oberwolfach Reports, DOI 10.4171/owr/2008/17, pp. 869--906. \url{https://doi.org/10.4171/owr/2008/17}.
\item Supplied reference: Douady, ``continuity of Julia sets,'' \url{https://doi.org/10.5802/xups.1994-03}.
\item The supplied \texttt{corrected\_verified} note reports checking content against the official report. The partially solved label and assessment checked 2026-08-21 are dated supplied triage requiring independent verification of exact dynamical hypotheses and status.
\end{itemize}
\end{document}
